\chapter{Libro, imprenta y circulación de ideas}
La imprenta modificó la producción y el consumo del conocimiento. Pettegree detalla cómo los primeros impresores adaptaron el mercado manuscrito a una economía de escala, generando catálogos, ferias del libro y redes de distribución pan-europeas.\cite{Pettegree2010} Estas redes multiplicaron la disponibilidad de textos clásicos y devocionales, facilitando la estandarización lingüística y la creación de públicos lectores.

\begin{figure}[H]
    \centering
    % Reemplazar por la imagen recortada correspondiente
    \includegraphics[width=0.75\textwidth]{FIGURES/prensa_gutenberg_placeholder}
    \caption{Prensa tipográfica y cajas de tipos móviles ("Pettegree, Andrew - The Book in the Renaissance (2010, Yale University Press) - libgen.li.epub", p. XX).}
    \label{fig:prensa-gutenberg}
\end{figure}

Las antologías y compilaciones, como la de Elmer y Webb, reflejan la maduración de un público que exigía textos fiables y comentados.\cite{Elmer2000} Bergin señala que esta cultura del libro impulsó reformas pedagógicas y la institucionalización de bibliotecas públicas y privadas, evidenciando el papel de la imprenta como tecnología de gobierno cultural.\cite{Bergin2015}
